% =====================================================================================
% PREMIER JET, 12 aout 2026. Justification de la dynamique du facteur systemique Y_{j,t},
% demandee par Caroline Hillairet au point d'etape : mouvement brownien pour le bruit de
% fond, processus a sauts pour les attaques ciblees.
%
% STATUT : NOTE, PAS ENCORE DU MEMOIRE. A relire avant integration au chapitre socle.
%
% CE QUE CETTE NOTE FAIT, ET SURTOUT CE QU'ELLE NE FAIT PAS. Elle ecrit la dynamique en
% temps continu dont le facteur actuel est l'increment normalise. Elle NE recalibre rien :
% la normalisation E[Y] = 0, Var(Y) = 1 est precisement ce qui laisse les probabilites
% marginales de defaillance, donc tous les SCR publies, a l'identique. La calibration est
% gelee depuis le 7 aout, et une note de redaction ne doit pas la deplacer.
%
% REVISION DU 12 AOUT, MEME JOUR. Caroline a demande si les sinistres arrivent regulierement
% ou groupes en fin d'exercice, en disant qu'il fallait regarder le PETIT t de Y_{j,t}. Le
% premier jet ne repondait pas : il ecrivait Y_{j,t} avec le meme indice que le terme de
% propagation D_{ik,t-1}, et sa recommandation (kappa infini) effacait la question au lieu de
% la traiter. Trois ajouts : le paragraphe des TROIS HORLOGES, la section GROUPEMENT qui
% rapatrie la mesure du script 08g, et une recommandation reecrite autour du nombre sans
% dimension theta = kappa x duree de grappe.
%
% DEUX OBJETS A NE JAMAIS CONFONDRE, et le projet s'est deja fait prendre par ce genre de
% confusion (quantile unitaire contre capital agrege, canaux contre piliers) :
%   - le facteur d'ETAT Y_{j,t}, objet de cette note, qui entre dans le probit de
%     defaillance de conformite du pilier j ;
%   - le melange de FREQUENCE de la binomiale negative, qui est une intensite commune
%     annuelle sur le comptage des sinistres.
% Les deux portent un aléa commun, ils ne vivent pas au meme etage du modele et ils ne
% s'additionnent pas.
% =====================================================================================

\documentclass[11pt,a4paper]{article}
\usepackage[utf8]{inputenc}
\usepackage[T1]{fontenc}
\usepackage[french]{babel}
\usepackage{amsmath,amssymb,amsthm}
\usepackage[margin=2.5cm]{geometry}
\usepackage{enumitem}
\usepackage{booktabs}   % \toprule, \midrule, \bottomrule des deux tables ajoutees le 12 aout

\newcommand{\Var}{\operatorname{Var}}
\newcommand{\E}{\mathbb{E}}
\newcommand{\Prob}{\mathbb{P}}

\title{La dynamique du facteur systémique $Y_{j,t}$ :\\ deux échelles d'aléa, et pourquoi
il en faut deux}
\author{Note de travail}
\date{12 août 2026}

\begin{document}
\maketitle

\section{Le point de départ, et ce qui manquait}

Le modèle retenu écrit la défaillance de conformité du pilier $j$ pour l'entité $i$ à la
date $t$ comme le franchissement d'un seuil par une variable latente :
\begin{equation}
\label{eq:latente}
  X_{ij,t} \;=\; \text{base}_j \;+\; \sum_{k\neq j} W_{kj}\,D_{ik,t-1}
  \;+\; s_j\,Y_{j,t} \;+\; \varepsilon_{ij,t},
  \qquad D_{ij,t} \;=\; \mathbf 1\{X_{ij,t} \geq 0\},
\end{equation}
où $\varepsilon_{ij,t}$ est l'aléa idiosyncratique, de variance $1$ par normalisation,
$s_j = \sqrt{\rho_j/(1-\rho_j)}$ la charge systémique du pilier et
$\text{base}_j = -K_j/\sqrt{1-\rho_j}$ son seuil.

Le terme $Y_{j,t}$ y était posé comme un facteur commun centré et réduit, sans que sa
\emph{dynamique} soit écrite. C'est une lacune réelle et pas une coquetterie : la loi de
$Y_{j,t}$ ne change pas les probabilités marginales de défaillance, mais elle gouverne
entièrement la façon dont plusieurs piliers tombent \emph{ensemble}. Or c'est exactement
l'objet du mémoire.

\section{Deux échelles d'aléa, et l'argument qui interdit de n'en garder qu'une}

\paragraph{Le bruit de fond.} Un système d'information subit en permanence un grand nombre
de petites atteintes : dérive de configuration, correctifs appliqués en retard,
hameçonnage ordinaire, incidents d'exploitation. Aucune ne déplace seule l'état de
conformité d'un pilier, et leur nombre sur un exercice est grand. Sous ces deux conditions,
la somme normalisée de ces contributions converge, par le théorème central limite sous la
condition de Lindeberg, vers une gaussienne. La modélisation naturelle de son accumulation
en temps continu est donc un \textbf{mouvement brownien}, et ce n'est pas un choix de
commodité : c'est la limite imposée par la structure de l'aléa, beaucoup d'événements dont
aucun ne domine.

\paragraph{Les attaques d'ampleur.} À l'autre extrémité se trouvent les campagnes qui
touchent simultanément de nombreuses entités par une vulnérabilité partagée. La condition
de Lindeberg y échoue précisément : \textbf{un seul événement domine la somme de l'année}.
Un mouvement brownien ne peut pas les représenter, non parce qu'il serait trop régulier en
trajectoire, mais parce que ses accroissements sont de variance proportionnelle au temps et
sans atome : il ne produit pas de déplacement d'ordre un en un temps d'ordre zéro. Ces
événements demandent une composante de \textbf{sauts}.

\paragraph{L'argument décisif est empirique, non stylistique.} On pourrait objecter qu'une
gaussienne de variance suffisante engloberait les deux régimes. La donnée l'interdit, et de
façon mesurable. Un facteur purement brownien conduit, sur les comptages annuels, à un
indice de dispersion égal à $1$. Les comptages observés donnent $1{,}19$ sur les cellules
firme-année et $2{,}24$ sur la série annuelle détendancée. La surdispersion est donc un
\emph{fait} et non un artefact, et elle est la signature d'un aléa commun à atomes : c'est
la composante de sauts qui la produit. Autrement dit, la seconde échelle n'est pas ajoutée
pour enrichir le modèle, elle est \textbf{identifiée par la surdispersion}.

\section{Trois horloges, et laquelle porte la normalisation}
\label{sec:horloges}

% AJOUT DU 12 AOUT, APRES UNE OBJECTION DE CAROLINE HILLAIRET AU POINT D'ETAPE : « est-ce que
% les sinistres arrivent tous les trois jours, ou tous sur les trois derniers jours de
% l'annee ? Il fallait regarder le petit t de Y_{j,t} ». L'objection est fondee et le premier
% jet de cette note ne lui repondait pas : il ecrivait Y_{j,t} avec le MEME indice t que le
% terme de propagation D_{ik,t-1}, si bien que la normalisation a variance un portait sur une
% echelle qui n'etait pas nommee. Il y a trois horloges, elles different de trois ordres de
% grandeur, et les confondre rend la formule ambigue.

L'écriture \eqref{eq:latente} porte un seul indice $t$, et c'est une ambiguïté qu'il faut
lever avant d'écrire quoi que ce soit sur $Y$. Le modèle fait intervenir \textbf{trois}
échelles de temps distinctes.

\begin{center}\small
\renewcommand{\arraystretch}{1.3}
\begin{tabular}{@{}l l l@{}}
\toprule
\textbf{Horloge} & \textbf{Ordre de grandeur} & \textbf{Ce qu'elle cadence} \\
\midrule
Pas de cascade & heures à quelques jours & le terme $\sum_k W_{kj} D_{ik,t-1}$ \\
Persistance du facteur, $1/\kappa_j$ & non identifiée & la mémoire de l'état systémique \\
Période comptable & un an & le quantile de charge agrégée, donc le capital \\
\bottomrule
\end{tabular}
\end{center}

\paragraph{La séparation des deux échelles extrêmes est vérifiée, pas supposée.} Le pas de
cascade se lit sur la donnée horodatée : l'excitation mesurée sur les dates au jour est
\emph{intra-journalière}, avec un pic à trois heures, et la vague de mai 2023 sur un
prestataire partagé s'étend sur cinq jours. La période comptable est l'année. Il y a donc
environ trois ordres de grandeur entre les deux, ce qui autorise le traitement usuel en
physique statistique comme en actuariat : \textbf{la cascade se résout instantanément à
l'échelle comptable}. C'est exactement ce que fait l'implémentation, qui tire pour chaque
sinistre l'ensemble des piliers atteints dans la loi de progéniture de la cascade, au lieu
de faire avancer un compteur d'années.

\paragraph{L'indice $t$ de \eqref{eq:latente} est donc le pas de cascade, et non l'année.}
Lue autrement, l'équation ferait franchir un lien de propagation par exercice, ce qui serait
absurde devant un pic d'excitation à trois heures. En revanche $Y_{j,t}$ est
\textbf{constant sur la durée d'une cascade}, puisque celle-ci est brève devant toute
persistance plausible du facteur : tous les piliers touchés par un même événement voient
la même valeur du facteur. Ce n'est pas une approximation de confort, c'est la traduction
formelle du mécanisme de choc commun.

\paragraph{Et la normalisation porte la période comptable.} C'est sur l'exercice que
l'incidence de fond $\Phi(-1{,}7) = 4{,}5\,\%$ et les charges systémiques $\rho_j$ ont été
calibrées. La contrainte $\E[Y] = 0$, $\Var[Y] = 1$ s'entend donc sur l'année, et c'est ce
qui rend l'enrichissement de la section suivante sans effet sur les probabilités publiées.

\section{L'écriture retenue}

Soit, sur l'espace filtré usuel, pour chaque pilier $j$ :
\begin{equation}
\label{eq:sde}
  \mathrm{d}Z_j(t) \;=\; -\kappa_j\,Z_j(t)\,\mathrm{d}t
  \;+\; \sigma_j\,\mathrm{d}B_j(t)
  \;+\; \mathrm{d}\!\left(\sum_{n\geq 1} \zeta^{(j)}_n\,
        \mathbf 1\{\tau_n \leq t\}\right)
  \;-\; \nu_j\,\mu_j\,\mathrm{d}t ,
\end{equation}
où
\begin{itemize}[itemsep=2pt,topsep=3pt]
\item $(B_1,\dots,B_5)$ est un mouvement brownien à cinq composantes, de matrice de
      corrélation instantanée $\mathrm{R}$ : c'est le \textbf{bruit de fond}, et
      $\mathrm{R}$ porte la co-occurrence \emph{ordinaire} entre piliers ;
\item $(\tau_n)$ est la suite des dates d'arrivée d'un processus ponctuel d'intensité
      $\nu_j$ et $\zeta^{(j)}_n$ les amplitudes associées, d'espérance $\mu_j$ : c'est la
      composante d'\textbf{attaques d'ampleur}. Le dernier terme est son compensateur, qui
      garantit $\E[Z_j(t)] = 0$ ;
\item $\kappa_j > 0$ est la vitesse de rappel. Elle n'est pas décorative : elle représente
      la remédiation, c'est-à-dire le fait qu'un écart de conformité se referme avec le
      temps. À $\kappa_j = 0$ on retrouve un facteur sans mémoire.
\end{itemize}

Le facteur qui entre dans \eqref{eq:latente} est alors l'état normalisé en fin de période :
\begin{equation}
\label{eq:norm}
  Y_{j,t} \;=\; \frac{Z_j(t)}{\varsigma_j},
  \qquad
  \varsigma_j^2 \;=\; \Var\!\left[Z_j(t)\right]
  \;=\; \underbrace{\frac{\sigma_j^2}{2\kappa_j}}_{\text{brownien}}
  \;+\; \underbrace{\frac{\nu_j\,\E\big[(\zeta^{(j)})^2\big]}{2\kappa_j}}_{\text{sauts}} ,
\end{equation}
la variance stationnaire d'un processus d'Ornstein-Uhlenbeck dirigé par une semimartingale
de Lévy se lisant composante par composante, les deux sources étant indépendantes.

\paragraph{Un paramètre nouveau, et il est interprétable.} La décomposition
\eqref{eq:norm} fait apparaître la \textbf{part de variance systémique portée par les
sauts} :
\begin{equation}
\label{eq:pi}
  \pi_j \;=\; \frac{\nu_j\,\E\big[(\zeta^{(j)})^2\big]}
                   {\sigma_j^2 + \nu_j\,\E\big[(\zeta^{(j)})^2\big]}
  \;\in\; [0,1].
\end{equation}
À $\pi_j = 0$ le pilier ne subit que du bruit de fond ; à $\pi_j = 1$ il ne subit que des
campagnes. C'est la grandeur que le point d'étape appelait de ses vœux, et elle se calibre
sur la surdispersion plutôt que de se poser.

\section{Ce que cette écriture change, et ce qu'elle ne change pas}

C'est le point le plus important de la note, et il doit être dit avant tout usage.

\paragraph{Elle ne change aucune probabilité marginale, donc aucun capital publié.} Le
probit \eqref{eq:latente} ne voit $Y_{j,t}$ qu'à travers sa loi centrée réduite. Tant que
l'on impose $\E[Y_{j,t}] = 0$ et $\Var[Y_{j,t}] = 1$, ce que fait \eqref{eq:norm} par
construction, la probabilité inconditionnelle de défaillance du pilier $j$ reste
$\Phi(\text{base}_j)$ et l'incidence de fond reste $\Phi(-1{,}7) = 4{,}5\,\%$. La
calibration est gelée : cette note l'explicite sans la toucher, et c'est délibéré.

\paragraph{Elle change la dépendance de queue, et c'est tout son objet.} À variance égale,
remplacer une gaussienne par un mélange brownien et sauts laisse les marges intactes et
\emph{augmente} la probabilité que plusieurs piliers franchissent leur seuil au cours du
même événement. Formellement, la composante de sauts introduit un atome commun dans la
filtration des cinq piliers, ce qu'une gaussienne multivariée ne produit jamais : le
coefficient de dépendance de queue supérieure d'une gaussienne de corrélation $r < 1$ est
nul, celui d'un mélange à choc commun est strictement positif. L'enrichissement se lit donc
là, et uniquement là.

\paragraph{Conséquence pour l'identification, et elle est favorable.} Puisque la composante
de sauts se manifeste par de la \emph{co-occurrence} et non par un déplacement des marges,
elle est identifiée par la même statistique qui identifie la partie symétrique $S$ de la
matrice de propagation. Elle ne touche pas la partie antisymétrique $A$, c'est-à-dire la
direction, qui reste non identifiée et bornée. La note ne dégrade donc pas le résultat
d'identification partielle : elle en respecte le partage.

\section{Les sinistres arrivent-ils groupés dans l'année ? Ce que la donnée dit}
\label{sec:groupement}

% CETTE SECTION REPOND A LA QUESTION POSEE AU POINT D'ETAPE, et elle le fait avec des chiffres
% qui existaient DEJA dans le projet, rangés sous l'etiquette « KPI de concentration P4 » et
% jamais relies au facteur systemique. C'est la meme classe de defaut que le facteur 3,34 ou le
% -53 % : la grandeur etait calculee, elle n'etait pas branchee la ou elle sert.
% RESERVE A LEVER : le pic d'excitation a trois heures est lu sur la figure O2, la sortie du
% script qui le calcule n'etant pas encore versionnee. A relancer sur le PC.

La question est empirique et elle a une réponse mesurée, sur $2\,922$ jours et $4\,614$
événements du secteur financier.

\begin{center}\small
\renewcommand{\arraystretch}{1.25}
\begin{tabular}{@{}l r@{}}
\toprule
Indice de Gini de la charge journalière & $0{,}69$ \\
Part des sinistres portée par le centile de jours le plus chargé & $16\,\%$ \\
\dots par les $2\,\%$ de jours les plus chargés & $21\,\%$ \\
\dots par les $5\,\%$ de jours les plus chargés & $34\,\%$ \\
Causes communes détectées & $49$, soit $6{,}1$ par an \\
Part des sinistres portée par ces causes communes & $20\,\%$ \\
Victimes par cause & médiane $11$, maximum $191$ \\
Indice de queue (Hill) du nombre de victimes par cause & $1{,}46$ \\
\bottomrule
\end{tabular}
\end{center}

\vspace{4pt}
\noindent Aucun des deux cas d'école n'est donc le bon. Les sinistres n'arrivent ni
régulièrement, ni tous dans une fenêtre unique de fin d'exercice : ils arrivent par
\textbf{une poignée de causes communes par an, à queue lourde sur le nombre de victimes}. La
vague de mai 2023 en donne la mesure, avec $388$ victimes en cinq jours consécutifs dont
$191$ sur une seule journée.

\paragraph{Un piège dans cette mesure, et il faut le nommer.} Les $191$ victimes d'une
journée sont $191$ \emph{entités distinctes} atteintes par une cause, et non une entité
subissant $191$ sinistres. Le Gini porte sur les comptes journaliers \emph{du marché} : ce
qu'il mesure est de la dépendance \textbf{en coupe}, pas une accumulation intra-annuelle
pour une entité donnée. Or le besoin de capital calculé est celui d'une entité. Les deux
questions sont distinctes et une seule est observée.

\paragraph{À l'échelle d'entité, la question est d'ailleurs presque vide.} Le mémoire publie
que $98{,}8\,\%$ des années y sont sans aucun incident matériel. Demander si les sinistres
d'une entité se regroupent dans l'année n'a donc guère de sens : il y en a au plus un. Le
regroupement observé est un phénomène de marché, et c'est à ce titre qu'il doit entrer dans
le modèle.

\paragraph{Et pour le quantile annuel, le motif d'arrivée est du second ordre.} La sévérité
est à queue régulièrement variable d'indice $1/\xi = 1{,}68$. La charge agrégée obéit alors
au \textbf{principe du grand saut unique} :
\begin{equation}
\label{eq:gsu}
  \Prob(S > x) \;\underset{x \to \infty}{\sim}\; \E[N]\;\Prob(X > x),
\end{equation}
c'est-à-dire que le quantile annuel est porté par \emph{un} sinistre et non par la somme de
plusieurs. Le premier terme de \eqref{eq:gsu} ne dépend que de $\E[N]$ et de la queue
marginale, deux grandeurs que le motif d'arrivée ne déplace pas. Que le sinistre qui porte
le quantile survienne en mars ou dans une grappe de décembre est donc sans effet au premier
ordre. Le mémoire chiffre déjà ce pont : à $\lambda = 21{,}56$, atteindre le quantile annuel
à $99{,}5\,\%$ demande un quantile par sinistre à $99{,}977\,\%$, et il reste un facteur
résiduel de $1{,}8$ qui est l'empreinte de la surdispersion et de la contagion.

\paragraph{Ce qui est en revanche du premier ordre, et c'est là que le modèle le met.} La
perte d'un sinistre n'est pas une sévérité isolée mais la somme des sévérités sur les
piliers que \emph{ce même} sinistre atteint :
\begin{equation}
\label{eq:perte-incident}
  L \;=\; \sum_{j \in \mathcal C} X_j ,
\end{equation}
où $\mathcal C$ est l'ensemble des piliers touchés, tiré dans la loi de progéniture de la
cascade. Le regroupement agit donc sur $\mathcal C$, pas sur les dates : il multiplie la
perte d'un événement unique. C'est exactement ce que portent la matrice de propagation et le
choc commun du pilier des tiers.

\paragraph{Conclusion, et elle est à déclarer plutôt qu'à supposer.} Le regroupement
temporel est réel, fort, et mesuré. Le modèle le représente \textbf{en coupe} et non
\textbf{dans le temps}, et ce choix est défendable pour trois raisons qui se cumulent : la
séparation des échelles du \S\ref{sec:horloges}, la vacuité de la question à l'échelle
d'entité, et le principe du grand saut unique \eqref{eq:gsu} à l'échelle du secteur. Ce qui
n'était pas acceptable était de le supposer en silence, ce que faisait le premier jet de
cette note.

\section{Sur l'auto-excitation, et pourquoi le processus de sauts n'est pas auto-excité}

La composante de sauts de \eqref{eq:sde} est un processus ponctuel dont l'intensité $\nu_j$
peut être constante, aléatoire, ou fonction du passé. Le troisième cas, l'auto-excitation,
a été testé et non retenu, et il importe de dire précisément contre quoi.

Le test n'a pas été mené contre un noyau de paille mais contre le noyau à retard
$\varphi(a) = \alpha\,a\,e^{-\beta a}$ et la logique en deux phases de la littérature
cyber-actuarielle de référence. Le résultat est net : l'excitation apparente est
\emph{intra-journalière}, avec un pic à trois heures, et le rapport de branchement passe de
$0{,}55$ à une valeur quasi nulle dès que les co-occurrences du même jour sont traitées
comme exogènes. Autrement dit, ce que le noyau auto-excité capte à l'échelle où le modèle
travaille n'est pas une contagion temporelle mais une \textbf{simultanéité}, laquelle est
déjà portée, et mieux, par la composante de choc commun.

Il faut souligner que le processus retenu \textbf{est} un processus à sauts, et même un
processus à sauts d'intensité aléatoire : la binomiale négative calibrée sur les comptages
est un mélange de Poisson, donc un processus doublement stochastique. Ce qui est écarté
n'est pas la nature de saut, c'est la \emph{rétroaction} du passé sur l'intensité. La
distinction n'est pas rhétorique : elle sépare ce que la donnée soutient de ce qu'elle ne
soutient pas, à la résolution disponible.

\section{Une variante, et la recommandation}

Deux écritures sont défendables et il faut choisir explicitement.

\begin{description}[itemsep=3pt,topsep=3pt]
\item[Forme à renouvellement.] $\kappa_j \to \infty$ dans \eqref{eq:sde} : le facteur est
      l'accroissement normalisé de la période, sans mémoire d'une période à l'autre. C'est
      la forme la plus sobre, elle reproduit exactement les chiffres publiés, et elle
      suffit à porter l'argument des deux échelles.
\item[Forme à mémoire.] $\kappa_j$ fini : l'état systémique persiste, ce qui est plus
      réaliste puisqu'une dégradation de conformité ne se referme pas en un exercice. Cette
      forme a un mérite propre, elle relie le facteur à la trajectoire multi-périodes des
      états de conformité déjà construite, dont $\kappa_j$ devient la vitesse de
      remédiation.
\end{description}

% RECOMMANDATION REECRITE LE 12 AOUT. La premiere version recommandait la forme a
% renouvellement sans un mot du regroupement intra-annuel, ce qui etait le plus mauvais choix
% possible : kappa infini signifie que le facteur se renouvelle instantanement, donc que le
% regroupement N'A PAR CONSTRUCTION AUCUN EFFET. Elle supposait donc la reponse a la question
% posee au point d'etape au lieu de la traiter. C'est une hypothese cachee, et le projet en a
% deja paye plusieurs.

\paragraph{Le nombre qui tranche entre les deux.} Ce n'est ni $\kappa_j$ seul ni la durée
d'une grappe seule, mais leur produit, seule quantité sans dimension du problème :
\begin{equation}
\label{eq:dimensionless}
  \theta_j \;=\; \kappa_j \times \text{(durée d'une grappe)} .
\end{equation}
À $\theta_j \ll 1$ le facteur ne bouge pas pendant la grappe : tous les sinistres qu'elle
contient voient la \emph{même} valeur de $Y_j$ et leurs sévérités sont fortement dépendantes.
À $\theta_j \gg 1$ le facteur s'est renouvelé entre deux sinistres de la grappe, qui
redeviennent quasi indépendants. La forme à renouvellement pose $\theta_j = \infty$.

\noindent \textbf{Recommandation.} Publier la \textbf{forme à renouvellement}, avec
\eqref{eq:sde} et \eqref{eq:norm} écrites en toutes lettres pour montrer de quoi elle est le
cas limite, et surtout avec le \S\ref{sec:groupement} qui \emph{déclare} ce que ce choix
suppose. Trois motifs, et ils ne se valent pas.

\begin{enumerate}[itemsep=3pt,topsep=3pt]
\item \textbf{Le motif de fond.} Le regroupement intra-annuel est du second ordre sur le
      quantile annuel, par \eqref{eq:gsu}, et du premier ordre sur la co-occurrence entre
      piliers, laquelle est déjà portée par la cascade et le choc commun via
      \eqref{eq:perte-incident}. Le mécanisme n'est donc pas ignoré, il est logé dans la
      coupe plutôt que dans le temps. C'est ce motif qui justifie le choix.
\item \textbf{Le motif d'identifiabilité.} $\kappa_j$ n'est pas identifiable sur la
      profondeur d'historique disponible, et publier un paramètre non identifiable au cœur
      du modèle coûte plus qu'il n'apporte.
\item \textbf{Le motif de discipline.} Calibrer $\kappa_j$ sur les durées de grappes
      observées serait une recalibration, non une correction d'erreur, et la calibration est
      gelée. Ce motif est le plus faible des trois : il dit quand le faire, pas s'il faut le
      faire.
\end{enumerate}

\paragraph{Ce que coûterait la forme à mémoire, si on la voulait quand même.} Il faudrait
estimer $\kappa_j$ sur la décroissance de l'autocorrélation des comptages sub-annuels,
rejouer le pipeline stochastique, et accepter que des centaines de nombres publiés bougent
sans qu'aucun déplacement soit imputable au seul $\kappa_j$. En contrepartie on gagnerait un
lien explicite entre le facteur et la trajectoire multi-périodes des états de conformité,
$\kappa_j$ y devenant la vitesse de remédiation. L'échange n'est pas absurde, mais il se
décide avant le gel et non après.

\paragraph{Une conséquence de rédaction, à tenir.} La forme retenue suppose
$\theta_j = \infty$. Cela doit être \textbf{écrit}, avec la mesure du
\S\ref{sec:groupement} en regard, et non laissé à la sagacité du lecteur. Une hypothèse qui
n'apparaît nulle part n'est pas une hypothèse : c'est un angle mort, et celui-ci a été trouvé
par une question posée en séance, pas par une relecture.

\end{document}
